Los experimentos fueron realizados en una máquina virtual con las siguientes características:

\begin{itemize}
\item \textbf{Hardware:} (máquina física)
    \begin{itemize}
    \item Procesador: Intel Core i5-10600K @ 4.10GHz
    \item Memoria Ram: 16GB DDR4 2666MHz
    \item GPU: NVIDIA GeForce RTX 3060ti
    \item Almacenamiento: 1TB SSD M.2 NVMe PCIe 3.0 
    \end{itemize}

\item \textbf{Detalles Máquina Virtual:}
    \begin{itemize}
    \item Sistema Operativo: Ubuntu 24.04.4 LTS
    \item Compilador: g++ 13.3.0
    \item Recursos Asignados a Máquina Virtual: 12 procesadores lógicos (hilos) y 7.7 GiB de memoria RAM.
    \end{itemize} 

\end{itemize} 

\textbf{Mediciones}

\begin{enumerate}
    \item \textbf{Tiempo de ejecución:} Se utilizó la biblioteca <chrono> de C++ para medir el tiempo transcurrido en ms durante la ejecución de cada algoritmo individualmente.
    Esto excluyendo operaciones de entrada/salida y compilación. Para la ejecución de los algoritmos de ordenamiento se realizaron 2 ejecuciones completas correspondientes al dataset
    generado. Los resultados de estas ejecuciones mostraron conscistencia entre sí (\cref{fig:comparacionEjecuciones}), lo que indica que el rendimiento de los algoritmos es estable y predecible bajo las condiciones de prueba establecidas.
    \item \textbf{Consumo de memoria:} Se empleó getrusage() para medir el consumo máximo residente (RSS) durante la ejecución de cada algoritmo. En el caso particular de los algoritmos de ordenamiento,
    se midió en dos ejecuciones, inicialmente solo medía el valor máximo RSS alcanzado por el proceso completo, no necesariamente la memoria extra usada por el algoritmo, por lo que 
se decidió medir el consumo de memoria antes y después de la ejecución del algoritmo, y restar ambos valores para obtener una estimación más precisa del consumo de memoria adicional.
    \item \textbf{Compilación:} Cada medición se estandarizó con la utilización de makefiles para compilar los algoritmos con las mismas opciones de optimización (-O3) y sin incluir 
    información de depuración (-g). Esto garantiza que las comparaciones entre algoritmos sean justas y consistentes.
    \item \textbf{Casos de prueba:} Se utilizaron sólo los datasets generados por los scripts correspondientes a cada sección.
\end{enumerate}



\textbf{Datasets}

Los datasets utilizados para los experimentos fueron generados utilizando scripts específicos para cada sección. 
Estos scripts se encargaron de crear conjuntos de datos con dos tipos de dominios: \texttt{D1}: \{0,1\}, \texttt{D7}: 
\{0,1,2,3,4,5,6,7,8,9\} en donde para algoritmos de ordenamiento la estructura de pruebas son arreglos de tipo \texttt{ascendente}, \texttt{descendente} y 
\texttt{aleatorio} dentro de los rangos de valores de $10^3$ hasta $10^7$ y para multiplicación de matrices de los tipos \texttt{dispersa}, 
\texttt{diagonal} y \texttt{densa} dentro de los rangos de valores de $2^2$ hasta $2^{10}$. Estos datasets proporcionaron una 
base sólida para evaluar el rendimiento de los algoritmos en diferentes escenarios y condiciones, permitiendo así obtener 
conclusiones más robustas sobre su comportamiento.